This chapter describes the construction of the synthetic training dataset, from
regulatory source extraction through archetype design to a fully operational
generator producing labeled transactions. In the absence of real transaction data
--- a common constraint in banking anomaly detection projects --- the quality of
the synthetic data determines the ceiling of model performance.

% ═══════════════════════════════════════════════════════════
\section{Source Extraction Methodology}
% ═══════════════════════════════════════════════════════════

The synthetic generator requires two classes of calibration data: parameters
describing the \textit{normal population} (how legitimate clients transact) and
patterns describing the \textit{flagged population} (what suspicious activity looks
like). No single source covers both, so a multi-source extraction strategy was
designed.

A structured extraction template was developed with seven trust dimensions
--- recency, authority, peer review, information volume, geographic relevance,
data type, and reproducibility --- each scored as Strong, Medium, or Weak. A
nine-dimension coverage matrix evaluates each source's contribution to the
generator's needs: amount distribution, operation type mix, transaction frequency,
inter-event time, temporal patterns, spatial/branch patterns, channel mix,
beneficiary/emitter known-rate, and operation sequence patterns. The last dimension
was added specifically for the LSTM, which requires cross-operation temporal
ordering to learn sequential anomalies.

Three sources were extracted using this template. A fourth (BCT circulars on
internal LAB/FT control) was identified but confirmed as restricted documents
not publicly available, documented as a known gap.

% ═══════════════════════════════════════════════════════════
\section{Source Extractions}
% ═══════════════════════════════════════════════════════════

\subsection{S1 --- PaySim}

PaySim is a synthetic mobile-money simulator based on real transaction logs from
an undisclosed African mobile-money provider. Its trust profile is 2~Strong (peer
review, reproducibility), 4~Medium, and 1~Weak (geographic relevance --- mobile
money in sub-Saharan Africa differs structurally from branch teller banking in
Tunisia).

PaySim's primary role is as a \textbf{beneficiary graph structure} source: the
dataset's agent-to-agent transfer network provides the best available model for
counterparty relationship patterns. Its secondary role is amount distribution
shape and day-of-month seasonality calibration. Operation mapping is direct for
three of four operations (\texttt{CASH\_OUT}~$\leftrightarrow$~retrait,
\texttt{CASH\_IN}~$\leftrightarrow$~versement,
\texttt{TRANSFER}~$\leftrightarrow$~virement). The largest gap is
\textbf{cheque}, which has no PaySim equivalent --- check-based operations are
absent from mobile-money ecosystems entirely.

The PaySim anomaly base rate of approximately 0.13\% was retained as the initial
calibration anchor for the generator's anomaly injection rate.

\subsection{S2 --- Amen Bank 2023 Annual Report}

The Amen Bank annual report provides \textbf{normal population calibration} ---
who the clients are and what the bank's infrastructure looks like. The document
has an internal two-tier trust structure: Tier~1 comprises audited financial
statements (signed by BDO Tunisie and La Générale d'Audit et Conseil), while
Tier~2 covers unaudited management narrative. This distinction guided which
figures were treated as hard constraints versus directional indicators.

Key extractions:

\begin{itemize}[leftmargin=2cm]
    \item \textbf{Client deposit composition:} 61\% individuals, 33\% private
          businesses, 6\% institutional. This directly calibrates the five
          archetype population proportions.
    \item \textbf{Branch network:} 150 branches across 9 geographic zones,
          establishing the branch pool size and zone structure.
    \item \textbf{Employee structure:} 1,139 total employees, 60.5\%
          network-assigned. Combined with the 150-branch figure, this yields
          the approximation of 2 tellers per branch used throughout the system.
    \item \textbf{Existing AI-AML system:} The report confirms Amen Bank already
          operates an AI-based AML system for bank-wide compliance reporting. This
          validates the project's positioning as a complementary system, not a
          replacement.
\end{itemize}

\subsection{S3 --- CTAF 2021 Annual Report}

The CTAF (Commission Tunisienne des Analyses Financières) annual report provides
\textbf{flagged population calibration} --- what suspicious operations look like
in the Tunisian regulatory context. Four typological case studies were extracted
using a structured 7-field template:

\begin{itemize}[leftmargin=2cm]
    \item \textbf{Case 1} (espèces-heavy structuring): Multiple versements just
          below the 10,000~DT reporting threshold, dispersed across branches.
          Informed the near-threshold ratio feature and the smurfing anomaly
          scenario design.

    \item \textbf{Case 2} (virement chain): Series of virements to unknown
          beneficiaries with round amounts. Informed first-time beneficiary
          detection and round-amount flagging logic in the Decision Service.

    \item \textbf{Case 3} (PPE exploitation): Politically exposed person used as
          a fund conduit. Validated the design decision to treat PPE as a
          structural flag layered on archetypes rather than a separate archetype.

    \item \textbf{Case 4} (remise de chèques laundering): Third-party cheques
          deposited in rapid succession. Informed cumulative amount tracking and
          cheque-specific anomaly patterns.
\end{itemize}

Six quantitative structural targets were also extracted: instrument co-occurrence
rates (espèces 68\%, virement 56\%, chèque 34\%), BBE (Billets de Banque
Étrangers) concentration around the 10,000~DT threshold (85\% of operations in
the 5,000--20,000~DT band), and predicate offence taxonomy (contrebande 20\%,
corruption 19\%, escroquerie 11\%).

An important calibration distinction: the Amen Bank report describes the normal
population (${>}99\%$ legitimate operations), while the CTAF report describes the
flagged population (${<}1\%$ that should trigger alerts). The synthetic generator
requires both distributions.

\subsection{Known Gap --- BCT Circulars}

BCT (Banque Centrale de Tunisie) circulars 2018-09, 2017-08, and 2021-05 would
have provided exhaustive regulatory threshold enumeration for internal LAB/FT
controls. These documents are restricted and not publicly available. The gap is
documented rather than papered over: the Decision Service's regulatory rules are
based on the thresholds extractable from the CTAF case studies and Law~n°41-2024,
not a complete BCT circular inventory.

% ═══════════════════════════════════════════════════════════
\section{Client Archetype Design}
% ═══════════════════════════════════════════════════════════

Five client archetypes were designed across six behavioral dimensions, calibrated
using the deposit composition data from S2 (Amen Bank annual report):

\begin{table}[H]
\centering
\caption{Client archetypes}
\label{tab:archetypes}
\begin{tabular}{@{}lcccl@{}}
\toprule
\textbf{Archetype} & \textbf{Population} & \textbf{Dominant Op.} &
\textbf{Monthly Vol.} & \textbf{Amount Range (DT)} \\
\midrule
Salaried      & 49\% & Retrait (50\%)    & $\sim$5.7  & 200--500 \\
Small Business & 33\% & Versement (55\%)  & $\sim$37.5 & 500--2,000 \\
Student       & 5\%  & Retrait (95\%)    & $\sim$7.2  & 40--80 \\
Retiree       & 7\%  & Retrait/Virement  & $\sim$3.2  & 200--500 \\
Big Business  & 6\%  & Virement (78\%)   & $\sim$63   & 2,000--20,000 \\
\bottomrule
\end{tabular}
\end{table}

Population proportions are derived from the 61\% individual / 33\% private
business / 6\% institutional split, with the individual segment subdivided
approximately 80\% salaried, 12\% retiree, 8\% student. The six behavioral
dimensions per archetype are: operation mix (probability distribution over
four operations), transaction frequency (mean and standard deviation per
operation per month), timing (preferred day of month), amount ranges (mean
and standard deviation per operation), counterparty loyalty (probability
of transacting with a known beneficiary/emitter), and branch loyalty
(probability of using the home branch).

All archetype parameters are externalized in a YAML configuration file,
enabling recalibration without code changes.

\subsection{Regulatory Discovery --- Law n°41-2024}

During archetype calibration, a regulatory discovery was made: Law~n°41-2024
(effective February 2025) caps individual cheques at 30,000~DT. BCT Q1~2026
data confirms a 24.9\% collapse in cheque usage volume following this law.
This required revising the big business archetype: remise de chèques frequency
dropped from approximately 12--15 per month to 2--4, with virement absorbing
the difference. A new anomaly scenario --- cheque structuring just under the
30,000~DT ceiling --- was added to the generator (see Section~\ref{sec:scenarios}).

% ═══════════════════════════════════════════════════════════
\section{Generator Architecture}
% ═══════════════════════════════════════════════════════════

The generator follows a \textbf{4-class, 2-phase} design that separates the
behavioral model from the anomaly injection mechanism:

\begin{enumerate}[leftmargin=2cm]
    \item \textbf{Archetype} (Layer 1): loads behavioral parameters from the
          YAML configuration file. Defines population-level priors --- what a
          ``typical salaried client'' looks like on average.

    \item \textbf{Client} (Layer 2): samples a fixed personality from the
          archetype prior via Gaussian sampling with clipping constraints.
          Each client's parameters (personal amount mean, personal frequency,
          personal preferred day) are drawn once at creation and remain fixed
          for the client's lifetime. Operation mix probabilities are sampled
          then normalized to sum to~1.

    \item \textbf{Generator} (orchestrator): executes in two phases.
          \textbf{Phase~1 (planning)} creates the client population, estimates
          total event volume, and assigns anomaly scenarios to selected clients
          with specific time windows and difficulty tiers.
          \textbf{Phase~2 (generation)} ticks through each simulated day,
          determines which clients transact (using a bell curve around their
          preferred day), generates normal or anomalous events depending on
          whether the client is in an active anomaly window, and produces the
          final chronologically ordered dataset.

    \item \textbf{Event} (inline): per-event noise application, payload
          construction based on operation type, branch selection respecting
          branch loyalty, and employee assignment from the visited branch's
          pool.
\end{enumerate}

The two-phase approach is essential for controlled anomaly injection. Hard-tier
scenarios like smurfing span 21 days and require sufficient normal history before
the anomaly window for the LSTM to have learned the client's baseline. Planning
the anomaly windows before generation ensures this temporal constraint is
respected.

\subsection{Branch-Linked Employee Assignment}

A critical data quality requirement: each employee is assigned to exactly one
branch (2 employees per branch, 300 total across 150 branches). The initial
implementation assigned employees from a global pool, allowing the same employee
to appear at multiple branches --- a physical impossibility in branch banking.
This was identified as a bug that contaminated the
\texttt{repeated\_client\_employee} anomaly scenario: the scenario requires a
specific client to visit a specific teller at a specific branch repeatedly, which
is meaningless if the teller can appear anywhere. After fixing employee assignment
to be branch-linked, the LSTM's AUC improved from 0.9430 to 0.9866 because
branch-linked employees produce cleaner sequential patterns.

% ═══════════════════════════════════════════════════════════
\section{Anomaly Scenario Design}
\label{sec:scenarios}
% ═══════════════════════════════════════════════════════════

Nine anomaly scenarios are organized across three difficulty tiers, reflecting
the detection challenge they present. Difficulty is defined by how many features
an anomaly modifies and over what time horizon.

\begin{table}[H]
\centering
\caption{Anomaly scenarios by difficulty tier}
\label{tab:anomaly-scenarios}
\begin{tabular}{@{}llp{6.5cm}l@{}}
\toprule
\textbf{Tier} & \textbf{Scenario} & \textbf{Description} & \textbf{Actor} \\
\midrule
\multirow{3}{*}{Easy (1 day)}
    & Amount spike      & Single transaction 5--10$\times$ client mean
                        & Client \\
    & Round amount      & Large round amount (5,000--20,000~DT)
                        & Client \\
    & Duplicate virement & Matching virement within short window
                        & Client \\
\midrule
\multirow{3}{*}{Medium (5 days)}
    & Frequency burst   & Transaction rate $>3\times$ daily average
                        & Client \\
    & Cumulative threshold & 24h/48h cumulative exceeds norm
                        & Client \\
    & Timing anomaly    & Transactions outside branch hours
                        & Client \\
\midrule
\multirow{3}{*}{Hard (21 days)}
    & Smurfing          & Versements in 8,000--9,800~DT range, multiple branches
                        & Client \\
    & Cheque structuring & Cheques near 30,000~DT ceiling (Law n°41-2024)
                        & Client \\
    & Repeated client-employee & Same client-teller pair $\geq$3 in 24h
                        & Cross \\
\bottomrule
\end{tabular}
\end{table}

Each scenario was validated for \textbf{feature-space detectability}: the injected
anomaly must produce a measurable change in at least one of the 30 enriched
features. Scenarios that modified only profile-level aggregates without affecting
per-event features were removed during a mid-project audit (reduced from 13 to 9
scenarios). For example, the original \texttt{behavioral\_drift} scenario gradually
shifted a client's operation mix over 21 days, but this change was absorbed by the
profile's 30-day distribution window and produced no detectable per-event signal.

The difficulty distribution follows a pyramid: 10\% easy, 30\% medium, 60\% hard,
reflecting the real-world observation that most suspicious activity is designed
to evade detection.

% ═══════════════════════════════════════════════════════════
\section{Data Validation}
% ═══════════════════════════════════════════════════════════

The generator produces four CSV files per run:

\begin{table}[H]
\centering
\caption{Generator output files}
\label{tab:generator-output}
\begin{tabular}{@{}llr@{}}
\toprule
\textbf{File} & \textbf{Content} & \textbf{Rows} \\
\midrule
\texttt{transactions.csv}      & All simulated events      & 243,904 \\
\texttt{clients\_master.csv}   & Client reference data     & 10,000 \\
\texttt{accounts.csv}          & Account reference data    & 10,000 \\
\texttt{employees\_master.csv} & Employee reference data   & 300 \\
\bottomrule
\end{tabular}
\end{table}

Key validation metrics on the final dataset:

\begin{table}[H]
\centering
\caption{Dataset validation}
\label{tab:data-validation}
\begin{tabular}{@{}ll@{}}
\toprule
\textbf{Metric} & \textbf{Value} \\
\midrule
Total events         & 243,904 \\
Anomaly rate         & 0.623\% (target: $\sim$0.13\%) \\
Anomaly scenarios    & 9 (all represented) \\
Unique clients       & 10,000 across 5 archetypes \\
Employees            & 300 (2 per branch, branch-linked) \\
Branches             & 150 \\
Cross-branch employee violations & 0 \\
Simulation period    & 180 days (January--June 2025) \\
Salaried retrait mean & $\sim$300~DT (consistent with Tunisian salary) \\
\bottomrule
\end{tabular}
\end{table}

The anomaly rate (0.623\%) is higher than the 0.13\% target from PaySim because
multi-day anomaly windows (5 days for medium, 21 days for hard) produce multiple
anomalous events per plan entry rather than one. This was accepted as a reasonable
trade-off: the rate remains well within the class imbalance range where
reconstruction-based and sequence-based models can learn effectively, and the
evaluation metrics (AUC-ROC) are threshold-independent.

\subsection{Oracle Profile Shortcut}

For model training, client personality parameters are exported directly from the
generator's in-memory objects to \texttt{clients\_master.csv}. This provides the
exact behavioral baseline needed for contrast feature computation without requiring
batch profile computation over the simulated transaction history. The oracle
shortcut produces the same 30 enriched features as the production streaming
pipeline, ensuring training-inference feature parity. In the production system,
these baselines are replaced by batch-computed profiles from real transaction
history --- the pipeline handles this substitution without code changes.